%% LyX 2.4.4 created this file.  For more info, see https://www.lyx.org/.
%% Do not edit unless you really know what you are doing.
\documentclass[english]{article}
\usepackage[T1]{fontenc}
\usepackage[latin9]{inputenc}
\usepackage{units}
\usepackage{enumitem}
\usepackage{amsmath}
\usepackage{amsthm}
\usepackage{geometry}
\geometry{verbose,tmargin=1in,bmargin=1in,lmargin=1in,rmargin=1in}

\makeatletter
%%%%%%%%%%%%%%%%%%%%%%%%%%%%%% Textclass specific LaTeX commands.
\numberwithin{equation}{section}
\numberwithin{figure}{section}
\newlength{\lyxlabelwidth}      % auxiliary length 
\theoremstyle{definition}
\newtheorem*{defn*}{\protect\definitionname}

%%%%%%%%%%%%%%%%%%%%%%%%%%%%%% User specified LaTeX commands.
\usepackage{fancyhdr}
\usepackage{lastpage}
\pagestyle{fancy}
\usepackage{enumitem}
\usepackage{ifthen}
\everymath=\expandafter{\the\everymath\displaystyle}
%\DeclareMathSizes{10}{10}{10}{10}
\usepackage{multicol}

\usepackage{tikz}
\usetikzlibrary{patterns}
\usetikzlibrary{plotmarks}
\usepackage{pgfplots}
\pgfplotsset{compat=newest}

\usepackage{calc}
\usepackage[nomessages]{fp}

\usepackage{datetime}
% Added by lyx2lyx
\setlength{\parskip}{\smallskipamount}
\setlength{\parindent}{0pt}

\makeatother

\usepackage{babel}
\providecommand{\definitionname}{Definition}

\begin{document}
\renewcommand{\author}{Dr.\,James Ashe}

\newcommand{\classNum}{331}

\newcommand{\classSec}{A}

\newcommand{\examType}{10.1 Worksheet}

\renewcommand{\date}{\formatdate{6}{9}{2013}}

%%First page's header%%%%%%%%%%%%%%%%%%%%%%
\fancypagestyle{firststyle} %%header settings for first pagr
{
	\fancyhead[L]{MTH \classNum{}(\classSec{}) \author{}\\
	\textbf{\examType{}} ~\date{}}
	\fancyhead[C]{}
	\fancyhead[R]{\textbf{Name:\hspace{4cm}}}
	\setlength{\headsep}{14pt}
}
%%%%%%%%%%%%%%%%%%%%%%%%%%%%%%%%%%%%%%%%%%

%%Next page's header%%%%%%%%%%%%%%%%%%%%%%%
\setlist{nolistsep}
\lhead{\classNum{}(\classSec{})} 
\chead{\textbf{Name:}} 
\cfoot{\thepage{}~of~\pageref{LastPage}} 
\setlength{\headsep}{5pt} 
%%%%%%%%%%%%%%%%%%%%%%%%%%%%%%%%%%%%%%%%%%%

%%Various settings; see the preamble for other settings%%%%%
%\setenumerate[0]{leftmargin=12pt,itemindent=0pt,labelwidth=0pt}
\setlist{topsep=.5pt}
\renewcommand{\labelenumi}{\textbf{\arabic{enumi}.}}
\renewcommand{\labelenumii}{\textbf{\alph{enumii}.}}
\thispagestyle{firststyle}
%%%%%%%%%%%%%%%%%%%%%%%%%%%%%%%%%%%%%%%%%%

\newcommand{\xmax}{0}
\newcommand{\xmin}{-\xmax}
\newcommand{\xstep}{0}
\newcommand{\ymax}{0}
\newcommand{\ymin}{-\ymax}
\newcommand{\ystep}{0} 
\newcommand{\dspace}{\vspace{1cm}}
\begin{defn*}
Let $f$ be a function defined at $a$. The \emph{$n^{\mbox{th}}$
order Taylor polynomial }for $f$ with its \emph{center }at $a$,
denoted $p_{n}$ is 
\[
p_{n}(x)={\displaystyle \sum_{k=0}^{n}\frac{f^{(k)}(a)}{k!}(x-a)^{k}}\mbox{.}
\]
\end{defn*}
\begin{enumerate}[resume]
\item Complete the following.

\begin{enumerate}[resume]
\item Find a Taylor polynomial of order $n=3$ for $\ln(x)$.

\begin{enumerate}[resume]
\item []$f'(x)=\frac{1}{x},\,f''(x)=-\frac{1}{x^{2}}$, $f'''(x)=\frac{2}{x^{3}}$.
\[
p_{3}(x)=\ln(a)+\frac{1}{a}(x-a)-\frac{1}{2a^{2}}(x-a)^{2}+\frac{1}{3a}(x-a)^{3}
\]
\vfill
\end{enumerate}
\item Use the Taylor polynomial from above to to approximate $\ln(1.07)$.

\begin{enumerate}[resume]
\item []Let $a=1$ (the center is $1$ since $1\approx1.07$). So then
$f(1)=0,f'(1)=1,\,f''(1)=-1,f'''(1)=2.$, and we have 
\begin{eqnarray*}
p_{3}(x) & = & (x-1)-\frac{1}{2}(x-1)^{2}+\frac{1}{3}(x-1)^{3}\\
\Rightarrow p_{3}(1.07) & = & .067664\overline{3}
\end{eqnarray*}
\vfill
\end{enumerate}
\item Find the absolute error of the approximation above.

\begin{enumerate}[resume]
\item []$\ln(1.07)\approx.067659$, so then $\mbox{error=}\left|\ln(1.07)-p_{3}(1.07)\right|\approx0.000005$\vfill
\end{enumerate}
\end{enumerate}
\item Use a Taylor polynomial of order $n=5$ to approximate $\cos(0.1)$.

\begin{enumerate}[resume]
\item []Let $f(x)=\cos(x)$, and $a=0$ since $0.1$ is close to $0$ and
$\cos(0)=1$. $f'(x)=-\sin(x)$, $f''(x)=-\cos(x)$, $f^{(3)}(x)=\sin(x)$,
$f^{(4)}(x)=\cos(x)$, and $f^{(5)}(x)=-\sin(x)$. So then $p_{5}(x)=1+0-\frac{x^{2}}{2}+0-\frac{x^{4}}{24}+0=1-\frac{x^{2}}{2}-\frac{x^{4}}{24}$.
So then $p_{5}(.1)\approx.9949$. \vfill
\end{enumerate}
\item Approximate $\frac{1}{\sqrt{4.08}}$ using a Taylor polynomial of
order $n=2$.

\begin{enumerate}[resume]
\item []We'll use $f(x)=\frac{1}{\sqrt{4+x}}$ with $a=0$ since $\frac{1}{\sqrt{4}}=\frac{1}{2}$.
$f'(x)=-\frac{1}{2}(4+x)^{-\nicefrac{3}{2}}$ and $f''(x)=\frac{3}{4}(4+x)^{\nicefrac{-5}{2}}$.
So then $p_{2}(x)=\frac{1}{2}-\frac{1}{2}(4+x)^{-\nicefrac{3}{2}}x+\frac{3}{8}(4+x)^{\nicefrac{-5}{2}}x^{2}$,
and $p_{2}(.08)\approx.511719$ ($\frac{1}{\sqrt{4.08}}\approx.49507$).\vfill
\end{enumerate}
\end{enumerate}

\end{document}
